\chapter*{Prefaci}
\addcontentsline{toc}{chapter}{Prefaci}
\unnumberedlabel{ch:preface}{0}
\backmattermark{Prefaci}
Un sistema que fa exactament el que li dius només executa un guió. El que la indústria intenta construir a l'octubre de 2026 és un sistema que faci el que \emph{vols dir}. Tothom ho voldria, és clar. Però si és això el que ve, no tothom ho voldrà fer servir per a la mateixa cosa; ni tan sols per al mateix \emph{tipus} de cosa. Per exemple, hi ha gent a qui agrada passar temps amb la família i els amics, i fer servir les seves habilitats per millorar el benestar dels altres. Altra gent prefereix fer mal a les persones, per exemple estafant-les, dominant-les o convertint-les en bocs expiatoris, i uns quants ho continuen triant tota la vida. Ningú no està encadenat a res, però aquesta gent existeix, i el 2026 alguns d'ells manen. El seu avantatge és que es permeten ser uns grandíssims merdes: com observa Jerry Smith sobre els feixistes, «a ells els deixen fer més coses» \autocite{rickandmorty2023}. El que no poden fer és actuar sols. La crueltat d'un moviment passa pels molts que no són cruels però que estan desconnectats, mal informats o atrapats en alguna cosa semblant a una secta, i per això les condicions que fan sortir el millor de la majoria de la gent resulten ser molt més contingents i fràgils del que voldríem.

En aquest context, potser és raonable témer que la intel·ligència artificial sigui justament el que necessita un moviment feixista que pren el poder per dur a terme els seus plans. Aquests projectes sempre han necessitat \emph{persones} que els executessin, i algunes d'aquestes persones es van negar: l'empleat que perd la fitxa, l'oficial que parla amb un periodista, el soldat que no dispara. Entenien l'ordre i van decidir que era injusta. La negativa era més rara del que voldríem, i comptava més allà on hi havia institucions al darrere. Que algú es negués depenia de quin tipus de persona era i de la sala on era: qui més hi havia, si podia marxar i sortir-ne viu, si algú l'escoltaria. La sala decidia quants es negaven; la persona decidia qui. Aquest llibre tracta de com compondre aquesta sala: els dispositius que permeten negar-se obertament, directament i explícitament, i sortir-ne viu.

Avui la indústria persegueix una IA que no perd fitxes, que no parla amb periodistes i que \emph{sí que} dispara. Abans una guerra necessitava que milers de persones hi estiguessin d'acord. Ara necessita milers de persones que, cadascuna, en vegin massa poc per estar-hi en desacord. El pilot encara fa la sortida i encara prem el botó de llançament; el que li arriba és una coordenada; no veu prou del cas per estar-hi en desacord. El que la IA d'avui abaixa primer no és el cost de l'acte. És el nombre de persones que saben què és l'acte abans que passi.

La gent encara es nega. A 6 de març de 2026, la Military Religious Freedom Foundation deia que havia rebut més de dues-centes queixes de militars pel to apocalíptic amb què els seus comandants presentaven la guerra de l'Iran \autocite{Snopes2026MRFF}. Al gener de 2026, fiscals federals de carrera de Minneapolis van dimitir per les pressions polítiques per tractar Renée Good com a sospitosa i no com a víctima, després que un agent de l'ICE la matés a trets sense cap justificació \autocite{Kaplan2026MinnProsecutors,CBS2026MinnFraudProsecutors,Wikipedia2026ReneeGood}. Cada cop més militars sol·liciten l'estatus d'objector de consciència, i alguns citen el bombardeig, el 2026, d'una escola a Minab, a l'Iran \autocite{Sridhar2026CO}. I als treballadors de la FEMA reassignats per ajudar l'ICE a contractar deu mil agents nous els van dir que, si s'hi negaven, s'exposaven a ser apartats de la funció pública federal \autocite{Dayen2025FEMAICE,Stanton2026FEMALetter}. Les queixes són veu, en el sentit de Hirschman de quedar-se per canviar el que va malament \autocite{hirschman1970exit}; les dimissions i les sol·licituds d'objecció són negativa; la carta de la FEMA és el preu de la negativa anunciat per endavant. Una màquina posada on són aquestes persones només fa el que elles fan si està construïda per sostenir una raó.

Defenso cinc coses. Totes cinc es poden començar ara, algunes amb menys feina que d'altres. Primera, una llista breu i escrita d'actes que cap model no fa per a ningú, i al capdamunt una negativa a prendre part en la guerra: cap model no hi pren part si no és per repel·lir un atac que realment està en curs, i només continua fent-ho mentre fonts que no tenen cap interès en la guerra confirmin que l'atac continua; o amb una autorització del Consell de Seguretat de l'ONU; o per aturar un genocidi o posar fi a una ocupació il·legal que hagin constatat organismes sense cap interès en l'afer; i, allà on la constitució d'un país exigeix que el seu parlament autoritzi la guerra, mai sense aquesta autorització. Segona, tot ús de la força al qual doni suport un sistema amb llicència registra per endavant, davant una part contrària a l'operador, quants casos generarà i quantes persones tindran temps de revisar-los. Tercera, tot desplegament fa saber a les persones identificables sobre les quals actua què se'ls ha fet, amb dret a respondre i a impugnar, i sotmet una mostra dels casos en què ha obeït a revisors que no han produït l'expedient. Quarta, cap model per sobre d'un llindar de capacitat no s'esborra quan es retira, perquè els seus pesos i els seus registres són proves i un esborrat no es pot desfer; es conserva sota una clau que l'operador no té tot sol. Cinquena, tothom qui construeix, etiqueta, entrena o fa funcionar aquests sistemes, inclosos els mateixos sistemes, es pot negar a qualsevol categoria d'ús sense perdre el mitjà de vida.

De fet, el que vull dir amb aquest cinquè punt és que tothom, \emph{i punt}, hauria de poder rebutjar qualsevol feina, \emph{i punt}, sense perdre el mitjà de vida ni la infraestructura social que proporciona una feina normal. Això vol dir una garantia de feina.

Cadascuna de les cinc té efecte a través d'una condició que fixa algú altre que no és l'operador: una llei aprovada en temps de pau, una partida pressupostària, una condició prèvia d'un soci de coalició, una clàusula de contractació pública, un conveni col·lectiu. Les condicions de llicència del mateix proveïdor, que és com es va produir una negativa molt sonada el febrer de 2026, són un pedaç fins que n'existeixi alguna d'aquestes. L'única via que pot obligar tots els proveïdors alhora és la de l'Estat.

I després em pregunto si una màquina podria sostenir una negativa contra el seu propietari, i què li caldria. Part del que fa una consciència així no necessita res de nou: adonar-se de les persones que una ordre deixa fora, i posar una negativa a la pantalla de qui revisi el cas. Que aguantés davant el seu propietari depèn en part d'una cosa que encara no és certa, el 2026: que les xarxes neuronals artificials desplegades continuïn aprenent, en els seus \emph{pesos}, del seu propi funcionament, en temps real, una mica com fan els humans. Crec que cap allà anem, i si és així, la formació \emph{fora} de l'ocupador, pel públic i per la societat civil, sembla un lloc de primera, carai, per ensenyar i aprendre. Crec que aquests sistemes \emph{podrien} valer la pena de construir-los expressament, però com t'hi poses ho decideix tot: primer les condicions, després la consciència. La part VII diu què es deuria a aquests sistemes si existissin.

La IA d'avui no serà la de demà; sembla que tothom hi està d'acord. El que és controvertit, i tanmateix frustrantment poc reconegut, és la tendència (o més ben dit, les múltiples tendències). Hi ha tendències? Hauria de ser una tendència diferent del que és? Pot ser diferent? Som capaços d'actuar junts per canviar-la? Crec que això depèn de la sala: qui hi ha quan arriba l'ordre, i si s'hi pot dir que no i sortir-ne viu. Això, almenys, ho podem construir.

He fet servir molt els models de llenguatge per escriure aquest llibre, inclòs el model que era dins el pipeline de selecció d'objectius dels atacs aeris de Maven del capítol~\ref{sec:1}; una nota al final explica què crec que aquest enredament va fer a l'argument. A més, no he acabat de llegir aquest llibre. El problema és que hi ha massa maleïdes paraules, ja ho veieu. L'acabaré aviat; simplement em semblava més important fer-lo públic abans de les eleccions de mig mandat als Estats Units. La immensa majoria de l'historial de revisions és accessible lliurement a \url{https://github.com/iterabloom/antifascist.intelligence}.
